\section{Related Work: Why Existing Approaches Fall Short}
\label{sec:related}

Table~\ref{tab:related} summarizes how five lines of work address the properties \voss{} needs. No line covers them
together; the paragraphs below state what each does and why it is insufficient here.

\begin{table}[t]
\caption{Coverage of the required properties by prior approaches (\ok{} addressed, \part{} partly, blank: not addressed).
The assessment is ours, from each approach's stated goals.}
\label{tab:related}
\centering\scriptsize
\setlength{\tabcolsep}{3pt}
\begin{tabular}{@{}p{2.45cm}cccccc@{}}
\toprule
 & \rotatebox{70}{Untrusted hosts} & \rotatebox{70}{Host churn} & \rotatebox{70}{Latency SLOs} & \rotatebox{70}{Plans trust and warm state} & \rotatebox{70}{Gray-failure detection} & \rotatebox{70}{Verifiable accounting} \\
\midrule
Volunteer computing \cite{anderson2004boinc}           & \ok & \ok &     &     & \part & \part \\
Serverless schedulers \cite{shahrad2020wild,fuerst2021faascache,ousterhout2013sparrow} &  &  & \ok & \part &  &  \\
Harvested-resource serverless \cite{zhang2021harvest}  &     & \ok & \ok & \part &  &  \\
Confidential serverless \cite{arnautov2016scone,wallet2025} & \ok &  & \ok &  &  &  \\
Trusted-component fault tolerance \cite{levin2009trinc,veronese2013minbft} & \ok &  &  &  &  &  \\
\midrule
\voss{}                                                & \ok & \ok & \ok & \ok & \ok & \ok \\
\bottomrule
\end{tabular}
\end{table}

\textbf{Volunteer computing.} BOINC-style systems harvest idle machines for scientific batch work
\cite{anderson2002seti,anderson2004boinc}. Because hosts are untrusted, they replicate tasks and accept a result when a quorum
agrees, and they use deadlines to tolerate hosts that disappear. Replication multiplies cost, presumes deterministic tasks without side effects, protects
integrity but not the confidentiality of inputs, and works on a time scale of hours. None of this fits invocations
that take milliseconds, may be non-deterministic, and carry private data. \voss{} instead removes the integrity and confidentiality
concern with attestation, so redundancy is needed only against omission, and then only for risky invocations.

\textbf{Serverless scheduling and keep-alive.} The cold-start literature decides when to keep sandboxes alive or pre-warm them
\cite{shahrad2020wild}, treats keep-alive as caching \cite{fuerst2021faascache}, or accelerates start-up with lighter
sandboxes and snapshots \cite{agache2020firecracker,oakes2018sock,ao2022faasnap,cadden2020seuss}. Low-latency schedulers
place short tasks by sampling \cite{ousterhout2013sparrow} and large-scale cluster managers pack tasks onto
fleets \cite{verma2015borg}. All assume trusted, effectively always-on machines. A warm sandbox is lost only through a
policy decision, never because its owner closed the lid, and the model has no resource corresponding to a provisioned
secret. \voss{} keeps the start-up-tier vocabulary but adds a per-session trust step and plans warm state under churn
and exposure constraints.

\textbf{Harvested and edge serverless.} Serverless platforms on harvested datacenter VMs handle evictions of
the underlying resources \cite{zhang2021harvest}, and edge schedulers place containers near data and users
\cite{rausch2021edge}. Their hosts are operated by the provider or a trusted party, so the failure model is eviction rather than silent omission by an
adversary, and no secret is entrusted to a party that might read it.

\textbf{Confidential serverless.} Trusted execution environments protect containers and functions from the host operating
system and the platform operator \cite{arnautov2016scone,costan2016sgx,amd2020snp}, and recent work targets serverless
functions in confidential VMs and containers \cite{wallet2025,segarra2024sesame}, and measurement studies quantify the cost of confidential VMs \cite{misono2025cvm}. These
designs address confidentiality on provider-managed server hardware. They do not address hosts whose sessions end
unpredictably, and they do not treat the establishment of trust as a scheduled resource with a cost, lifetime and exposure.
\voss{} is substrate-agnostic (Section~\ref{sec:substrates}): the same scheduler runs over a TPM-attested appliance on consumer
hardware or over confidential VMs, with the substrate profile supplying establishment costs and assurance.

\textbf{Trusted components and failure detection.} Small trusted components (a TPM-like counter or log) restrict what a faulty node can do
and reduce replication needed for Byzantine agreement \cite{levin2009trinc,veronese2013minbft}. We use the same idea
differently: attestation does not make agreement cheaper; it makes the faulty node \emph{non-Byzantine}, which lets us replace
agreement among volunteers by a provider-run arbiter (Section~\ref{sec:succession}). For detection, heartbeat and accrual
detectors \cite{hayashibara2004phi} infer liveness, which is not progress; gray failures
\cite{huang2017gray} motivate observing end-to-end progress, which \voss{} does at the trusted router and the attested scheduler.

\textbf{Hedging and robust provisioning.} Hedged requests reduce tail latency by duplicating slow work
\cite{dean2013tail}. \voss{} applies hedging \emph{in proportion to estimated risk}, on a different failure domain, and
bounds the extra work by a budget. For the planner we use cardinality-constrained (budgeted) robust optimization
\cite{bertsimas2004price}, which fits the question ``is demand still covered if any $f$ hosts defect together?''
without assuming a probability model for adversaries.
